% Memoria autónoma para incorporar después al informe. Requiere amsmath, booktabs y biblatex.
\section{Actualización del eje transversal: M8 y acoplamiento de contacto}
Se estudia una varilla de acero al carbono zincado M8 por1,25, comparada con M6 por1
bajo condiciones equivalentes. No se modifica el STEP ni se certifica el grado comprado.
Se emplean núcleos de referencia $6,466$ y $4,773\,\mathrm{mm}$, según la tabla de TR Fastenings\parencite{tr_rosca_m6_2026},
con $I=\pi d^4/64$. No se utiliza el área resistente a tracción como sección de flexión.
El módulo de referencia $E=210\,\mathrm{GPa}$ corresponde a acero al carbono, según SCI/BCSA\parencite{sci_acero_elastico_2026};
el recubrimiento zincado no identifica grado ni límite elástico. Se contrasta $E=200/220\,\mathrm{GPa}$.
Los pesos usan densidades de referencia y volúmenes CAD; no son masas medidas.

\subsection{Geometría y restricciones}
Se conserva la longitudCAD de $160\,\mathrm{mm}$. Ambos empotramientos se sitúan en los centros
de los taladros laterales de cámara, separados $108,411756\,\mathrm{mm}$, imponiendo traslación y giro
nulos en ambos planos. Los extremos físicos restantes están en voladizo. No se confunde esta luz
con la longitud total ni se atribuye rigidez infinita comprobada al ABS.
Los centros de rodillos y cuatro rodamientos se conservan del ensamble auditado.
El diámetro de varilla M8 es una propuesta analítica: los taladros de cámara del CAD siguen siendo
de $6\,\mathrm{mm}$. Su incompatibilidad de montaje debe resolverse sin editar aquí el diseño.
La coincidencia nominal M8 con bore608 de $8\,\mathrm{mm}$ no verifica asiento roscado,
ajuste, alineación ni apriete. El alojamiento de rodillo de $24\,\mathrm{mm}$ sigue sin ajustarse
al exterior608 de $22\,\mathrm{mm}$.

\subsection{Algoritmo recuperado de EP2, modelo3C}
Se reutilizan la integral de superposición variable y la viga Euler--Bernoulli de EP2\parencite{equipo_entrega2_2026}.
Se parametrizan su geometría, presión y condiciones de borde, sin modificar los parámetros globales comerciales.
Para cada uno de los1440estados se actualiza
\begin{align}
g_j(r)&=g_{0j}+v_{cj}+\alpha_{rj}(r-r_{cj}),\\
L_c(r)&=\sqrt{D_r\max(h-g_j(r),0)-\max(h-g_j(r),0)^2},\\
F_j&=P\sum_i A_{ij}(g_j),\qquad T=\mu_{\mathrm{ef}}\sum_{j,i}r_iP A_{ij},\\
\mathbf K\mathbf u&=\mathbf f(\mathbf g),\qquad
\mathbf s^{k+1}=(1-\lambda)\mathbf s^k+\lambda\mathbf s_{\mathrm{calc}}^k.
\end{align}
Se usan $P=2,258670672\,\mathrm{MPa}$, $\mu_{\mathrm{ef}}=0,1$ y las capas supuestas1/2/4mm,
sin seleccionar una condición nominal física. La presión es teórica y uniforme, no medida ni constitutiva.
Las fuerzas se reparten en cada par608 por equilibrio de fuerza y momento usando el centroide de área.
Los desplazamientos de sus centros determinan la traslación e inclinación del rodillo rígido.
Se transforma la pendiente global al eje radial propio de cada rodillo; no se reutilizan signos de posición comerciales.
El torque de la matriz no se aplica como torsión adicional al eje transversal.

La relajación de referencia es $\lambda=0,5$, con máximo100iteraciones por estado.
Se exige residuo de compatibilidad y cambio de perfil de holgura inferiores a $10^{-9}\,\mathrm{m}$,
y cambio relativo de las fuerzas por rodillo inferior a $10^{-5}$, con denominador mínimo1N,
como en EP2. El contacto se recalcula al final y se reverifica contra la deformación exportada.
Un estado no convergido se registra como tal y no se usa para iniciar el siguiente;
las holguras negativas se rechazan sin recortarlas para forzar solución.

\subsection{Resultados y límites}
\begin{table}[ht]\centering
\begin{tabular}{lrrrrr}\toprule
Varilla/capa&$F_{\max}$ [N]&$T_{\max}$ [N m]&$u_{\max}$ [mm]&$\sigma_{b,\max}$ [MPa]&Iter.\\\midrule
M6x1/1mm & 350.98 & 1.0776 & 0.1868 & 285.86 & 16\\
M6x1/2mm & 603.12 & 1.8315 & 0.3257 & 491.26 & 18\\
M6x1/4mm & 894.97 & 2.7188 & 0.5134 & 762.70 & 19\\
M8x1.25/1mm & 377.64 & 1.1317 & 0.0604 & 124.69 & 17\\
M8x1.25/2mm & 629.86 & 1.9025 & 0.1025 & 208.50 & 18\\
M8x1.25/4mm & 918.00 & 2.7804 & 0.1569 & 315.36 & 18\\
\bottomrule\end{tabular}
\caption{Casos acoplados con fijaciones hipotéticas idénticas. Tensiones nominales elásticas, sin concentraciones.}
\end{table}
\begin{table}[ht]\centering\begin{tabular}{lrrrr}\toprule
Varilla/capa&$g_{1,\min}$&$g_{1,\max}$&$g_{2,\min}$&$g_{2,\max}$\\\midrule
M6x1/1mm & 0.5220 & 0.6740 & 0.5226 & 0.6744\\
M6x1/2mm & 0.5378 & 0.8033 & 0.5390 & 0.8041\\
M6x1/4mm & 0.5716 & 0.9777 & 0.5736 & 0.9789\\
M8x1.25/1mm & 0.5070 & 0.5563 & 0.5072 & 0.5564\\
M8x1.25/2mm & 0.5117 & 0.5954 & 0.5121 & 0.5957\\
M8x1.25/4mm & 0.5218 & 0.6460 & 0.5224 & 0.6463\\
\bottomrule\end{tabular}\caption{Holguras convergidas en mm; envolventes angulares y radiales por rodillo.}\end{table}
La condición geométrica de holgura positiva y menor que la capa se verifica separadamente
de la aceptación funcional. No existe una banda de holgura operativa admisible documentada.
Mantener holgura positiva no significa mantener exactamente los $0,5\,\mathrm{mm}$ iniciales.
Tampoco la convergencia matemática acredita el régimen elástico real: faltan grado, límite elástico,
raíz/tolerancia de rosca, concentradores, rigidez de montaje y adaptación608.
No se presenta factor de seguridad, fatiga ni aprobación del montaje. La matriz y la carcasa
se consideran rígidas para esta iteración, no como resultado FEM.
